Large language models are increasingly used to simulate diverse personas, yet it remains unclear whether some personas require extensive context conditioning to emerge. We investigate whether \emph{context-heavy personas}---personas that can only be reliably elicited after conditioning on thousands of tokens of in-context examples---exist in frontier language models. Using \gptmini on 25 diverse personas from the Anthropic Model-Written Evals benchmark, we systematically vary the number of in-context demonstrations $K$ from 0 to 200 and measure persona elicitation accuracy. We find \emph{no evidence} for context-heavy personas: all 25 personas achieve $\geq$95\% action consistency at $K{=}0$ (zero-shot), with additional context providing negligible improvement (mean gain $<$0.4\%). A second experiment measuring open-ended generation distinctiveness via LLM-as-judge evaluation finds only marginal improvement with more context (mean gain +0.18/10, $p{=}0.038$). We term this the \emph{shallow persona hypothesis}: for frontier models, all tested personas are specifiable with short prompts, suggesting the accessible persona space is bounded by what can be concisely described rather than by context-length-dependent activation thresholds. Our results indicate that the practical persona space of frontier LLMs, while combinatorially vast, is shallow in depth---no extensive contextual buildup is required.
